% page 115 of the facsimile (image p-114 of the scan)
% block 172.7 x 246.3 mm ; left margin 29.3 mm
\pgc{7.620mm}{0.000mm}{
\l{\cel{57}l05}
\l{}
\l{}
\l{\cel{1}\sou{deviacar}. (\sou{netrans.,}\cel{1}de). Eskartesar de la voyo rekta ; eskartesar}
\l{\cel{4}de la direciono normala. - EFIS.}
\l{}
\l{\cel{1}\sou{devizo}. Formulo quan onu skribis super, sub od apud la skudo en la}
\l{\cel{4}blazono, olim. - DEFIRS.}
\l{}
\l{devoco.- Zelo por religio, por la praktikado religiala. - DEFIS..}
\l{}
\l{\cel{1}\sou{deocrar}. (\sou{trans.}) I. (\sou{aludante animalo}). Saturar su per sua rapt-}
\l{\cel{4}ajo. - II. Manjar avide. - EFIS.}
\l{}
\l{\cel{1}\sou{dvvot}a.Konsakrita a servar od a sorgar la interesti di ulu. -DEFI.}
\l{}
\l{dextra. Situita ye la latero opozita ad olta ube trovesas la pintoo}
\l{\cel{4}di le homo-kordio. - EFIS.}
\l{}
\l{\cel{1}\sou{dextrino}. (\sou{kemio}).Substanco analoga a gumo, di qua la kompozanti}
\l{\cel{4}kemiala esas sama kam olti di amilo, e quan onu uzas por apreta,}
\l{\cel{3}bluizar.\cel{2}Simbolo kemiala :\cel{2}(C6 H10 O5) n. - DEFIRS..}
\l{}
\l{\cel{1}\sou{dezeraa}. Qua ne esas sovaja (od arida), ma ne-habitata. - EFIS.}
\l{}
\l{\cel{1}\sou{dezinenco}. Fino-literi di vorto, qui expresas la flexioni. - DFIS.}
\l{}
\l{dezirar. (trans., por).\cel{2}Tendencar a (kozo quan onu volas posedar,,}
\l{\cel{3}od ago quan onu volas efektigar - segun quante ta voli esass}
\l{\cel{4}satisfacebla). - EFIS.}
\l{}
\l{\cel{1}\sou{dezolar}. (\sou{trans.}) Frapar (ulu) per aflikto troa, pro qua judikeble}
\l{\cel{4}omno mankas a lu. - EFIS.}
\l{}
\l{\cel{1}di. (\sou{gram.}) Prepoziciono qua expresas la relato di persono o kozo:}
\l{\cel{4}l. kun termino a qua ta persono o kozo apartenas, o de qua lu}
\l{\cel{4}dependas ; 2. kun termino submisita a lua efiko.}
\l{}
\l{\cel{1}\sou{dio}. Tempo-quanto qua fluas (inter du pasi da Suno ye la sama meri-}
\l{\cel{4}diano, de dimezo a dimezo), e qua varias segun la epoko di la}
\l{\cel{3}yaro (t.e. segun la situeso di Tero sur ekliptiko). - EFIS..}
\l{}
\l{\sou{diabaso}. (\sou{geol}.) Roko eruptala, qua havas la sama kompozanto mine-}
\l{\cel{4}ralogiala kam diorito, e strukturo ek intrikajo de prismi aplas-}
\l{\cel{4}tita, feldspata, aglomerita en plaji de piroxeno. - DEF.}
\l{}
\l{\cel{1}\sou{diabeto}. (\sou{patol}.)Morbo karakterizata da ek-sekreco abundanta di}
\l{\cel{3}urino qua kontenas plu o min multa elementi sukra - e da deperiso}
\l{\cel{4}gradoza. - DEFIRS.}
\l{}
\l{\cel{1}\sou{diablo}. I. (en\cel{1}\sou{la religio Kristaaa}). La demono di lo mala, Satanas,}
\l{\cel{4}princo di la anjeli rebela. - II. Persono o kozo quan onu kom-}
\l{\cel{4}paras a diablo. - DEFIRS.}
\l{}
\l{\cel{1}diado. I. Asemblajo de du principi, de du elementi. - II. (\sou{Filoz}.)}
\l{\cel{4}Duo ek la principo unesma (monado) e la principo inferiora, qua}
\l{\cel{4}devenas de ica.}
}
